\chapter{CONCLUSION AND RECOMMENDATIONS}

\ChapterSection{Conclusion}

Krishi Vaidya was developed as an integrated crop-health support system that combines computer
vision, vision-language modelling, backend orchestration, and a mobile application into one
project. The work demonstrates that a final-year B.Sc. CSIT project can move beyond a single
model experiment and produce a connected software system with trained model components,
application services, mobile workflows, structured data storage, and validation artifacts.

The YOLO component establishes the object-detection layer of the system. It was trained on a
curated plant-part detection dataset containing 30,981 images and 147,467 object annotations
across eleven agricultural object classes. The validation and result analysis show that the model
achieved usable object-detection behavior for the intended plant-part localization task, with the
best recorded training metrics reaching precision 0.71717, recall 0.74868, mAP50 0.76163, and
mAP50-95 0.60191. The deployment study further examined threshold behavior and quantized
variants, which is important because the target system includes mobile and edge-oriented
constraints.

The VLM component establishes the disease-interpretation layer. The project prepared a structured
vision-language dataset, validated the dataset pipeline, fine-tuned a Qwen-based VLM using LoRA
and quantized training techniques, and evaluated prediction behavior through recorded validation
outputs. The VLM work is technically significant because it supports diagnosis-oriented text and
structured output rather than only fixed-label classification. At the same time, the current VLM
results reveal an important limitation: the recorded prediction metrics and parse-error behavior
show that the VLM pipeline is implemented and measurable, but not yet reliable enough to be
treated as a final diagnostic authority without further improvement.

The backend component provides the application control layer. It implements authentication,
domain routes, data models, repository abstraction, scan handling, diagnosis orchestration,
advisory generation, validation, and test coverage for the scan path. The backend is not only a
transport layer between the mobile app and the models; it is the component that turns model
output into a structured application response and maintains system records.

The mobile application provides the user-facing layer. It implements routing, screen structure,
provider and repository patterns, local SQLite persistence, offline and synchronization support,
crop and task management, image-based scan workflow, backend diagnosis communication, model asset
support, localization, and farmer-facing output. This makes the project usable as an application
prototype rather than only a collection of scripts.

The integrated workflow connects these components into a diagnosis path: the mobile app captures
or selects an image, the backend receives and validates the scan request, the VLM service handles
model inference, the backend transforms the result into an advisory response, and the mobile app
stores and displays the outcome. The integration has been verified from source structure and
documented data contracts, while live end-to-end runtime evidence remains a required final
validation step.

Overall, Krishi Vaidya satisfies the core objective of building an integrated agricultural
diagnosis-support prototype. Its strongest current evidence lies in the completed YOLO artifacts,
the complete VLM data and training pipeline, the implemented backend, the implemented mobile app,
and the documented integration architecture. Its most important remaining work is not basic
implementation, but reliability improvement, executed end-to-end testing, mobile runtime
evidence, field evaluation, and expert validation of diagnosis and advisory content.

\ChapterSection{Major Contributions}

The major contributions of this project are as follows:

\begin{enumerate}[label=\arabic*.]
\item A curated YOLO object-detection dataset and trained detector for agricultural plant-part
localization, including validation metrics, threshold analysis, quantization comparison, and
error analysis.
\item A structured VLM dataset and fine-tuning pipeline for crop disease interpretation using
Qwen-based vision-language modelling, LoRA adaptation, quantized training, prediction parsing,
and validation reporting.
\item A backend API that organizes authentication, crop records, scan records, diagnosis
services, advisory logic, repository abstraction, validation middleware, and scan-related tests.
\item A React Native/Expo mobile application that supports farmer-facing scan workflow, local
records, offline persistence, synchronization structure, crop and task management, localization,
and backend diagnosis communication.
\item A documented integration workflow that connects the mobile app, backend, VLM inference
service, and advisory response into a single diagnosis-support path.
\item A manuscript-ready planning and evidence system that maps model outputs, figures, tables,
architecture diagrams, limitations, and remaining validation requirements.
\end{enumerate}

\ChapterSection{Recommendations}

Krishi Vaidya should be continued as an evidence-first system. The next work should not only add
features, but should strengthen the reliability of the existing components. The recommendations
are grouped into dataset and model improvements, deployment improvements, and mobile/user
evaluation improvements.

\ChapterSubsection{Dataset and model improvements}

The YOLO dataset should be expanded with more balanced class coverage. The current object
distribution is useful for training, but some classes have much fewer annotations than others.
Additional images for under-represented classes such as maize ear, rice grain cluster,
cauliflower head, and cabbage head would improve the detector's ability to generalize across
plant structures. More field images should also be collected under varied lighting, occlusion,
growth stages, camera quality, and background conditions.

The VLM dataset should be improved with stricter output supervision. The current validation shows
that parse errors are a major limitation. Future dataset preparation should enforce consistent
JSON-style response templates, include negative and uncertain cases, and add examples where the
correct behavior is to decline overconfident diagnosis. The VLM should also be evaluated with
expert-reviewed labels and advisory text rather than relying only on automated parsing metrics.

The diagnosis pipeline should separate disease classification confidence from advisory
confidence. A model may identify a likely disease while still lacking enough context to recommend
chemical or agronomic intervention. Future versions should include confidence calibration,
uncertainty labels, and escalation guidance such as ``consult an agricultural expert'' when the
visual evidence is weak.

\ChapterSubsection{System deployment improvements}

The backend and VLM services should be tested under a live end-to-end deployment. The manuscript
currently documents source-verified integration, but final deployment should include request logs,
response examples, latency measurements, failure-mode testing, authentication checks, and
readiness behavior. These results should be archived and added to the appendix before final field
deployment claims are made.

The backend test suite should be expanded beyond scan-focused tests. Additional tests should cover
authentication, crops, notes, tasks, advisories, diagnosis history, validation failures, file
upload limits, and authorization boundaries. This is necessary because the backend is responsible
for data integrity and user-specific records.

The model-serving boundary should be hardened. The VLM service should expose documented health,
readiness, timeout, and error responses. The backend should handle VLM unavailability without
crashing the user workflow. Where possible, model outputs should be schema-validated before being
stored or returned to the mobile app.

\ChapterSubsection{Mobile app and user evaluation improvements}

The mobile app should be tested on real Android devices and, if available, different device
classes. The final evidence should include screenshots, scan workflow recordings, local storage
behavior, offline-to-online synchronization, response latency, error messages, and farmer-facing
readability checks. These tests are necessary because mobile usability cannot be fully proven from
source code inspection alone.

The application should also be evaluated with intended users. Farmer-facing language, disease
names, severity explanations, and advisory instructions should be checked with agricultural
experts and local users. If Nepali-language output is used, translation quality must be evaluated
carefully because unclear advice can be more harmful than no advice.

Future versions should include stronger safety boundaries. The mobile app should avoid presenting
model output as a guaranteed diagnosis. It should show uncertainty, recommend expert consultation
for severe or unclear cases, and avoid unsupported pesticide prescriptions unless the advisory
content has been reviewed by qualified agricultural personnel.

\ChapterSection{Future Enhancements}

The following enhancements are recommended for future development:

\begin{enumerate}[label=\arabic*.]
\item Add a complete end-to-end test run from mobile image submission to backend response, VLM
inference, advisory generation, local storage, and result display.
\item Improve VLM structured-output reliability by training with stricter response schemas and
adding parser-aware validation during evaluation.
\item Add expert-reviewed disease and advisory labels for the VLM validation set.
\item Expand the YOLO dataset with more balanced class representation and more field-condition
images.
\item Add final mobile screenshots, device-test evidence, and offline synchronization
demonstrations to the appendix.
\item Add backend API documentation or Swagger/OpenAPI-style endpoint summaries for easier
maintenance.
\item Add deployment documentation for environment variables, model-service startup, backend
startup, mobile configuration, and production readiness checks.
\item Conduct a small controlled user evaluation with students, farmers, or agricultural
technicians to assess usability and clarity of advisory output.
\item Add privacy and data-protection review for uploaded crop images and user records before
real-world deployment.
\item Extend Nepali-language support for disease explanations and farmer-facing recommendations
after expert review.
\end{enumerate}
